\documentclass{article}
\usepackage[utf8]{inputenc}
\usepackage{amsmath}
\usepackage{geometry}
\geometry{margin=2cm}
\begin{document}

\section*{Questão 93 — ENEM 2024 — Ciências da Natureza}

\subsection*{Tema: Comunicação Sonora em Anfíbios}
\textbf{Interpretação da questão:} Esta questão exige compreender qual é a função do som estridente emitido por anfíbios da espécie \emph{Eleutherodactylus jahnstonei} encontrados em São Paulo.

\subsubsection*{Termos-chave}
\begin{description}
  \item[Emitir:] Produzir ou enviar sinais (neste caso, sons) para comunicação.
  \item[Função:] O propósito ou papel do som emitido pelos animais.
  \item[Anfíbios Anuros:] Grupo de animais que inclui sapos e rãs, caracterizados principalmente pelo som que produzem.
\end{description}

\subsection*{Análise e Estratégia}
Para responder, é necessário analisar o comportamento dos anuros, que utilizam sons para comunicação, especialmente com finalidade reprodutiva. O enunciado indica um som constante e estridente que causa incômodo, comum em chamadas para atração de parceiros. Examinando as alternativas, deve-se identificar aquela que condiz com o uso biológico do som por esses animais.

\subsection*{Mini revisão conceitual}
\begin{tabular}{|p{5cm}|p{5cm}|p{5cm}|}
\hline
\textbf{Conceito} & \textbf{Ideia-chave} & \textbf{Aplicação} \\ \hline
Comunicação animal & Animais emitem sons para diversos objetivos, incluindo atração sexual, alerta e interação social. & Anfíbios anuros usam canto para atrair fêmeas durante a reprodução. \\ \hline
Função do canto em anuros & O canto é uma estratégia para encontrar parceiros e garantir a reprodução da espécie. & O som estridente tem como função avisar a presença e atrair fêmeas. \\ \hline
Ecologia de espécies exóticas & Animais fora de seu habitat natural podem continuar exibindo comportamentos típicos. & A espécie do Caribe presente no Brooklin mantém comportamento de canto reprodutivo. \\ \hline
\end{tabular}

\subsection*{Aspectos teóricos relevantes}
\begin{itemize}
  \item Os anfíbios da ordem Anura, como sapos e rãs, possuem uma comunicação sonora altamente desenvolvida, cuja principal função é a atração de fêmeas durante o período reprodutivo, através do canto, comportamento instintivo que ocorre independente do local.
  \item O canto dos anuros é emitido principalmente pelos machos para marcar território e sinalizar disponibilidade sexual, fator crucial para o sucesso reprodutivo.
  \item Mesmo em ambientes urbanos, algumas espécies exóticas mantêm seus comportamentos naturais, como a emissão de sons altos para comunicação sexual, que pode causar incômodo a humanos.
\end{itemize}

\subsection*{Resolução da Questão}
\begin{itemize}
  \item O enunciado informa que o som estridente ocorre em área residencial da cidade, causando desconforto.
  \item É sabido que o canto dos anuros é uma forma de sinal acústico para comunicação sexual, especialmente atração de fêmeas para reprodução.
  \item Analisando as alternativas, a correta é aquela que identifica a função do canto como atração sexual.
  \item Alternativas que atribuem outras funções ao som, como estresse, alerta ambiental, superpopulação ou atração de presas, não condizem com conhecimento biológico.
  \item Logo, a alternativa correta é a \textbf{E}.
\end{itemize}

\subsection*{Pulo do gato}
\emph{Lembre-se:} o canto dos anuros serve para atrair fêmeas e garantir a reprodução da espécie, e não para comunicar problemas ambientais ou sociais.

\subsection*{Análise das alternativas}
\begin{itemize}
  \item \textbf{A:} Som não indica estresse dos animais – é o canto natural utilizado para reprodução.
  \item \textbf{B:} Anuros não utilizam canto para alertar sobre poluição ou perigo ambiental.
  \item \textbf{C:} Canto não indica superpopulação, mas busca por parceiros sexuais.
  \item \textbf{D:} Sons não atraem presas; atraem fêmeas para acasalamento.
  \item \textbf{E:} Alternativa correta – som serve para atrair fêmeas durante o acasalamento.
\end{itemize}

\subsection*{Código da questão}
CN\_BIOLOGIA\_COMPORTAMENTO\_001

\end{document}